\section{Engineering and Reproducibility}

The project targets Python 3.10--3.12 and separates configuration, data, networks, policy,
environment construction, and evaluation into testable modules. YAML defines the experiment, and
the resolved JSON configuration is copied into every run directory. Checkpoints contain online and
EMA weights, optimizer and scheduler state, EMA state, current iteration, best metrics, and the
training configuration. This is why the 15k continuation is mathematically the same scheduled run
rather than an informal warm start.

The measured environment is summarized in Table~\ref{tab:software}. The Pod used the mutable
\texttt{maniskill/base:latest} image, so exact package versions are recorded to remove ambiguity.

\begin{table}[t]
    \centering
    \caption{Software and hardware for the completed Push-T run.}
    \label{tab:software}
    \footnotesize
    \begin{tabular}{@{}ll@{}}
        \toprule
        Component & Version / device \\
        \midrule
        Python / ManiSkill / SAPIEN & 3.12.13 / 3.0.1 / 3.0.3 \\
        PyTorch / Diffusers & 2.13.0 / 0.40.0 \\
        Gymnasium / NumPy / h5py & 1.3.0 / 2.5.2 / 3.16.0 \\
        GPU / driver-reported CUDA & NVIDIA A40 / 13.0 \\
        \bottomrule
    \end{tabular}
\end{table}

The evaluator writes one \texttt{metrics.json} per seed, and a separate aggregation command writes
\texttt{summary.json}. Complete evidence includes the aligned HDF5 and metadata hashes, config,
checkpoint lineage, both training logs, TensorBoard events, five-second GPU CSV, seed metrics,
aggregate summary, and videos. Large artifacts are backed up outside ordinary Git history; source,
configs, tests, launchers, and concise result summaries remain in Git.

Implementation safeguards include:
\begin{itemize}
    \item CPU vector environments expose episode metrics directly, whereas the GPU wrapper may
    place them under \texttt{final\_info}; the evaluator supports both representations.
    \item Underscore-prefixed fields such as \texttt{\_success\_once} are Gymnasium validity masks,
    not metrics. Aggregation uses the corresponding non-prefixed value.
    \item All evaluation environments use \texttt{reconfiguration\_freq=1}; otherwise cached scene
    configurations can make nominally different episodes share initialization state.
    \item RGB replay aborts on controller, parallelism, observation-count, or source-code mismatch
    instead of silently producing a different dataset.
    \item Unit tests cover configuration validation, temporal sequence construction, policy shapes,
    CPU/GPU metric extraction, three-seed aggregation, replay alignment, and dataset provenance.
\end{itemize}

The local verification suite contains 20 passing T-I tests. The repository exposes equivalent
Colab, HCE, and Runpod entry points, but the minimal reproduction path is the stage-based Runpod
launcher in Appendix~\ref{sec:reproduction}.
